% Part: first-order-logic
% Chapter: proof-systems
% Section: axiomatic-deduction

\documentclass[../../../include/open-logic-section]{subfiles}

\begin{document}

\iftag{FOL}
      {\olfileid{fol}{prf}{axd}}
      {\olfileid{pl}{prf}{axd}}

\olsection{સ્વયંસિદ્ધિમૂલક \usetoken{P}{derivation}}

સ્વયંસિદ્ધિમૂલક !!{derivation}s સૌથી જૂનાં અને સૌથી સરળ તાર્કિક
!!{derivation} તંત્રો છે. તેમાંની !!{derivation}s કેવળ !!{sentence}sની
શ્રેણીઓ હોય છે. !!{sentence}sની કોઈ શ્રેણીને યોગ્ય !!{derivation}
ગણવા માટે તેમાંનું દરેક !!{sentence}~$!A$ નીચેની શરતોમાંથી કોઈ એક
સંતોષતું હોવું જોઈએ:
\begin{enumerate}
\item $!A$ સ્વયંસિદ્ધિ છે, અથવા
\item $!A$ એ આપેલા ગણ~$\Gamma$નું !!a{element} છે અને તે ગણ
  !!{sentence}sનો છે, અથવા
\item $!A$ને અનુમાનનો કોઈ નિયમ પ્રમાણિત કરે છે.
\end{enumerate}
સ્વયંસિદ્ધિ બનવા માટે $!A$નું સ્વરૂપ કેટલાક નિશ્ચિત !!{sentence}
પ્રરૂપોમાંથી કોઈ એક જેવું હોવું જોઈએ. પ્રથમ-ક્રમના તર્કશાસ્ત્ર માટે
સંતોષકારક (યથાર્થ અને પૂર્ણ) !!{derivation} તંત્ર આપતા સ્વયંસિદ્ધિ-
પ્રરૂપોના ઘણા ગણ છે. કેટલાક પ્રરૂપો તેઓ જે સંયોજકોને નિયંત્રિત કરે છે
તે મુજબ ગોઠવેલા હોય છે; દાખલા તરીકે,
\[
!A \lif (!B \lif !A) \qquad !B \lif (!B \lor !C) \qquad (!B \land !C) \lif !B
\]
પ્રરૂપો અનુક્રમે $\lif$, $\lor$ અને~$\land$ને નિયંત્રિત કરતી સામાન્ય
સ્વયંસિદ્ધિઓ છે. કેટલાંક સ્વયંસિદ્ધિ-તંત્રો સ્વયંસિદ્ધિઓની સંખ્યા
ઓછામાં ઓછી રાખવાનું લક્ષ્ય ધરાવે છે. કયા સંયોજકો મૂળભૂત ગણાય છે તેના
પર આધાર રાખીને માત્ર એક સ્વયંસિદ્ધિ ધરાવતાં તંત્રો પણ મળી શકે છે.

અનુમાનનો નિયમ એક એવું શરતી વિધાન છે જે !!a{sentence}ને
!!a{derivation}માં પ્રમાણિત કરવા માટેની પૂરતી શરત આપે છે. મોડસ પોનેન્સ
આવો એક બહુ સામાન્ય નિયમ છે: તે કહે છે કે જો $!A$ અને $!A \lif !B$
પહેલેથી પ્રમાણિત હોય, તો $!B$ પ્રમાણિત છે. તેનો અર્થ એ કે
!!a{derivation}માં !!{sentence}~$!B$ ધરાવતી પંક્તિ પ્રમાણિત છે, જો
$!A$ અને $!A \lif !B$ (કોઈ !!{sentence}~$!A$ માટે) બંને~$!B$ની પહેલાં
તે !!{derivation}માં આવે છે.

સ્વયંસિદ્ધિમૂલક !!{derivation}s પર આધારિત $\Proves$ સંબંધ આ પ્રમાણે
વ્યાખ્યાયિત થાય છે: $\Gamma \Proves !A$ જો અને માત્ર જો એવું
!!a{derivation} હોય જેનું છેલ્લું સૂત્ર !!{sentence}~$!A$ હોય (અને
ઉપરની શરત~(2)થી પ્રમાણિત થયેલાં !!{sentence}s જે !!{derivation}માં
આવે છે તેમનો ગણ $\Gamma$ લેવાય). જો $!A$ને એવી !!a{derivation} હોય
જેમાં $\Gamma$ ખાલી હોય—એટલે કે દરેક !!{sentence} જે !!{derivation}માં
આવે છે તે કાં તો~(1) વડે અથવા~(3) વડે પ્રમાણિત હોય—તો $!A$ પ્રમેય છે.
દાખલા તરીકે, નીચેની
!!a{derivation} બતાવે છે કે $\Proves !A  \lif (!B \lif (!B \lor !A))$:
\begin{derivation}
  1. & $!B \lif (!B \lor !A)$ \\
  2. & $(!B \lif (!B \lor !A)) \lif (!A  \lif (!B \lif (!B \lor !A)))$\\
  3. & $!A  \lif (!B \lif (!B \lor !A))$
\end{derivation}
પંક્તિ~1 પરનું !!{sentence} સ્વયંસિદ્ધિ $!A \lif (!A \lor !B)$ના
સ્વરૂપનું છે ($!A$ અને $!B$ની ભૂમિકાઓ ઊલટાવીને). પંક્તિ~2 પરનું
વાક્ય સ્વયંસિદ્ધિ $!A \lif (!B \lif !A)$ના સ્વરૂપનું છે. તેથી બંને
પંક્તિઓ પ્રમાણિત છે. પંક્તિ~3ને મોડસ પોનેન્સ પ્રમાણિત કરે છે: જો તેને
ટૂંકમાં $!D$ લખીએ, તો પંક્તિ~2નું સ્વરૂપ $!C \lif !D$ છે, જ્યાં $!C$
એ $!B \lif (!B \lor !A)$, એટલે કે પંક્તિ~1 છે.

જો $\Gamma \Proves \lfalse$, તો ગણ $\Gamma$ અસુસંગત છે. પૂર્ણ
સ્વયંસિદ્ધિ-તંત્ર કોઈપણ~$!A$ માટે $\lfalse \lif !A$ પણ સાબિત કરશે;
તેથી જો $\Gamma$ અસુસંગત હોય, તો કોઈપણ~$!A$ માટે
$\Gamma \Proves !A$.

તર્કશાસ્ત્ર માટેનાં સ્વયંસિદ્ધિમૂલક !!{derivation}sનાં તંત્રો સૌપ્રથમ
ગોટ્લોબ ફ્રેગે પોતાની 1879ની કૃતિ \emph{Begriffsschrift}માં આપ્યાં
હતાં; આ કારણે તેને ઘણી વાર આધુનિક તર્કશાસ્ત્રની પ્રથમ કૃતિ ગણવામાં
આવે છે. આ તંત્રોને આલ્ફ્રેડ નૉર્થ વ્હાઇટહેડ અને બર્ટ્રાન્ડ રસેલના
\emph{Principia Mathematica}માં તથા 1920ના દાયકામાં ડેવિડ હિલ્બર્ટ અને
તેમના વિદ્યાર્થીઓએ વધુ પરિપૂર્ણ બનાવ્યાં. તેથી તેમને ઘણી વાર
``ફ્રેગ-તંત્રો'' અથવા ``હિલ્બર્ટ-તંત્રો'' કહે છે. આવાં તંત્રો બહુ
બહુમુખી છે, કારણ કે કોઈ તર્કશાસ્ત્ર માટે સ્વયંસિદ્ધિમૂલક તંત્ર શોધવું
ઘણી વાર સહેલું હોય છે. !!{derivation}sની રચના બહુ સરળ હોય છે અને તેમાં
માત્ર એક કે બે અનુમાન-નિયમો હોય છે, તેથી તેમના \emph{વિશે} પરિણામો
સાબિત કરવાં પણ પ્રમાણમાં સહેલાં છે. પરંતુ વ્યવહારમાં તેમનો ઉપયોગ કરવો
ઘણો અઘરો છે, એટલે કે સાબિતીઓ શોધવી અને લખવી મુશ્કેલ છે.

\end{document}
